The stated identity is **false as written**. Here is a concrete counterexample.

Use the standard convention that a residue is the coefficient of the \(h^{-1}\) term in the Laurent expansion. ([mathworld.wolfram.com](https://mathworld.wolfram.com/ComplexResidue.html?utm_source=openai))

Take
\[
l_1=1,\qquad l_2=0,\qquad l_3=-3.
\]
Then
\[
l_1+l_2+l_3=-2,\qquad l_1\ge l_2\ge l_3,\qquad l_1\ge 1,
\]
so the hypotheses are satisfied. Also \(l_2\le 0\), hence by definition
\[
R_B=0.
\]

Now choose
\[
\lambda_1=1,\qquad \lambda_2=2,\qquad \lambda_3=-\lambda_1-\lambda_2=-3.
\]
Since \(l_1=1\), we have
\[
\Phi_{l_1}(z)=\Phi_1(z)=z.
\]
For these values,
\[
q_A(h)=\frac{(h+2)(h-1)(h-2)}{h(h+1)(h-4)}.
\]
Near \(h=0\),
\[
q_A(h)=\frac{4+O(h)}{-4h+O(h^2)}=-\frac1h+O(1).
\]
Also
\[
P_A(h)=(h-1)^2(h-1)^{-3}(h+1)(h-4)^4(h-2)^{-4},
\]
so
\[
P_A(0)=(-1)^{-1}\cdot 1\cdot (-4)^4\cdot (-2)^{-4}=-16.
\]
Therefore
\[
P_A(h)\Phi_1(q_A(h))
=
P_A(h)q_A(h)
=
\left(-16+O(h)\right)\left(-\frac1h+O(1)\right)
=
\frac{16}{h}+O(1).
\]
Hence
\[
R_A=16,
\qquad
R_A+R_B=16.
\]

Now compute the claimed right-hand side. We have
\[
\bar l_1=1,\qquad \bar l_2=0,\qquad \bar l_3=2,
\]
and therefore
\[
\sigma(\bar l_1)=1,\qquad \sigma(\bar l_2)=0,\qquad \sigma(\bar l_3)=3.
\]
Thus
\[
Q
=
1\cdot 1^{-2}
+
0
+
3\cdot (-3)^{-2}
+
0
+
0
+
(1)(-3)(1^{-1})(-3)^{-1}.
\]
So
\[
Q
=
1+\frac{3}{9}+1
=
\frac73.
\]
The claimed right-hand side equals
\[
\frac18\left(
\lambda_1^{l_1+1}\lambda_2^{l_2+1}\lambda_3^{l_3+1}
-
\lambda_1\lambda_2\lambda_3 Q
\right)
=
\frac18\left(
1^2\cdot 2^1\cdot (-3)^{-2}
-
(1)(2)(-3)\cdot \frac73
\right).
\]
That is
\[
\frac18\left(\frac29+14\right)
=
\frac18\cdot \frac{128}{9}
=
\frac{16}{9}.
\]

But
\[
R_A+R_B=16\neq \frac{16}{9}.
\]

Therefore the proposed identity cannot be proved under the definitions given; at least one of the definitions or the claimed formula must contain a typo or missing normalization.